% 論文節の日本語訳（読み物 PDF）用 preamble
% lualatex + jlreq
\documentclass[12pt,paper=a4]{jlreq}
\usepackage{luatexja-fontspec}
\usepackage{amsmath,amssymb,bm}
\usepackage{graphicx}
\usepackage{hyperref}
\usepackage{enumitem}
\usepackage{booktabs}
\usepackage{microtype}
\usepackage{xcolor}
\usepackage{caption}
\usepackage{tcolorbox}
\tcbuselibrary{breakable,skins}

\hypersetup{
  colorlinks=true,
  linkcolor=blue!60!black,
  urlcolor=blue!60!black,
  citecolor=blue!60!black,
}

\graphicspath{{images/}{./}}

\setlength{\parindent}{1em}
\setlength{\parskip}{0.55em}
\setlist[itemize]{leftmargin=1.6em,itemsep=0.25em,topsep=0.35em}
\setlist[enumerate]{leftmargin=1.8em,itemsep=0.25em,topsep=0.35em}

\captionsetup{
  font=small,
  labelfont=bf,
  skip=0.4em,
  justification=raggedright,
  singlelinecheck=false,
}

\newcommand{\vx}{\bm{x}}
\newcommand{\vu}{\bm{u}}
\newcommand{\veps}{\bm{\epsilon}}
\newcommand{\vw}{\bm{w}}

\DeclareMathOperator*{\argmax}{arg\,max}
\DeclareMathOperator*{\argmin}{arg\,min}
\DeclareMathOperator{\Tr}{Tr}
\DeclareMathOperator{\diag}{diag}

% 命題・定理ボックス（ページ跨ぎ時も切れ目を抑える）
\newtcolorbox{propositionbox}[1][]{%
  enhanced,
  breakable,
  colback=blue!4!white,
  colframe=blue!45!black,
  coltitle=white,
  fonttitle=\bfseries\small,
  title={#1},
  attach boxed title to top left={yshift=-2mm, xshift=3mm},
  boxed title style={
    colback=blue!45!black,
    colframe=blue!45!black,
    boxrule=0pt,
    arc=1pt,
  },
  boxrule=0.5pt,
  arc=1.5pt,
  left=5pt, right=5pt, top=8pt, bottom=5pt,
  before skip=0.9em,
  after skip=0.9em,
  pad at break=3mm,
}

% 図はフロートせず本文位置に置く（分断を防ぐ）
\newcommand{\translationfigure}[2]{%
  \begin{center}
    \includegraphics[width=0.88\linewidth]{#1}\\[0.35em]
    {\small #2\par}
  \end{center}
  \vspace{0.4em}%
}

\begin{document}
{\color{gray}\footnotesize 出典: §2.2\par}
\vspace{0.2em}

\section*{2.2 変分的視点：DDPM}


Denoising Diffusion Probabilistic Models（DDPMs）（Sohl-Dickstein et al., 2015; Ho et al., 2020）は拡散モデリングの要石である。概念的には、VAE や HVAE と同様、変分枠組みの中で動作する。しかし DDPM は、先行モデルが直面した課題のいくつかに対処する巧妙なひねりを導入する。

その核心において、DDPM は二つの異なる確率過程を含む：

\begin{itemize}
\item \textbf{前向き過程（固定エンコーダ）}: この過程は、遷移核 $p(x_i\mid x_{i-1})$ を通じて複数ステップにわたってガウスノイズを注入し、データを徐々に劣化させる。データは等方的ガウス分布へと発展し、実質的に純ノイズとなる。すなわちエンコーダは固定であり、学習されない。
\item \textbf{逆デノイジング過程（学習可能なデコーダ）}: ここでは、ニューラルネットワークがパラメータ化された分布 $p_\phi(x_{i-1}\mid x_i)$ を通じてノイズ劣化を逆転するよう学習する。純ノイズから始め、この過程は反復的にデノイズして現実的なサンプルを生成する。決定的に、個々の各デノイジングステップは、VAE がしばしば試みるように完全なサンプルを一から生成するよりも扱いやすい課題である。
\end{itemize}


エンコーダを固定し、学習を漸進的な生成軌道に集中させることで、DDPM は顕著な安定性と表現力を達成する。

\translationfigure{fig-2-3.jpg}{図2.3: DDPM の図示。データにガウスノイズを徐々に加える固定前向き過程（灰色）と、新しいサンプルを生成するためにステップごとにデノイズする学習された逆過程からなる。}


本節では DDPM に焦点を当て、より一般的で柔軟な枠組みを提示する第4.5節へ広い議論は延期する。

\subsection*{2.2.1 前向き過程（固定エンコーダ）}


DDPM では、前向き過程はエンコーダとして機能する固定・非学習の操作である。複数ステップにわたってノイズを加えることで元データを徐々に劣化させ、最終的に単純な事前分布 $p_{\mathrm{prior}} := \mathcal{N}(0,I)$ へと変換する。この変換は図2.3の前向き連鎖、あるいは図2.4に示される通りである。

このステップごとの劣化を形式化しよう：

\translationfigure{fig-2-4.jpg}{図2.4: DDPM 前向き過程の図示。ガウスノイズが段階的に加えられ、データサンプルを純ノイズへと劣化する。}


\textbf{固定ガウス遷移。} 前向き過程の各ステップは固定ガウス遷移核\footnote{この定式化は見かけ上異なるように見えるかもしれないが、元の DDPM 遷移核と数学的に等価である。}で支配される：

\[
p(x_i\mid x_{i-1}) := \mathcal{N}\bigl(x_i;\, \sqrt{1-\beta_i^2}\, x_{i-1},\, \beta_i^2 I\bigr).
\]


ここで過程は $x_0$ から始まり、これは実データ分布 $p_{\mathrm{data}}$ から抽出されたサンプルを表す。系列 $\{\beta_i\}_{i=1}^{L}$ は事前に定められた単調増加のノイズスケジュールを表し、各 $\beta_i \in (0,1)$ はステップ $i$ で注入されるガウスノイズの分散を制御する。便宜上 $\alpha_i := \sqrt{1-\beta_i^2}$ と定義する。この数学的定義は、次の直感的な反復更新と正確に等価である：

\[
x_i = \alpha_i x_{i-1} + \beta_i \epsilon_i,
\]


ここで $\epsilon_i \sim \mathcal{N}(0,I)$ は独立同分布である。すなわち各ステップ $i$ で、前の状態 $x_{i-1}$ を $\alpha_i$ で縮小し、$\beta_i$ でスケールされた制御された量のガウスノイズを加える。

\textbf{摂動核と事前分布。} 遷移核を再帰的に適用することで、元データ $x_0$ が与えられたときのステップ $i$ におけるノイズ付きサンプルの分布の閉形式表現が得られる：

\[
p_i(x_i\mid x_0) = \mathcal{N}\bigl(x_i;\, \bar{\alpha}_i x_0,\, (1-\bar{\alpha}_i^2)I\bigr),
\]


ここで

\[
\bar{\alpha}_i := \prod_{k=1}^{i}\sqrt{1-\beta_k^2} = \prod_{k=1}^{i}\alpha_k.
\]


これは、$x$ から $x_i$ を直接サンプリングできることを意味する：

\[
x_i = \bar{\alpha}_i x_0 + \sqrt{1-\bar{\alpha}_i^2}\,\epsilon, \qquad \epsilon\sim\mathcal{N}(0,I). \quad (2.2.1)
\]


ノイズスケジュール $\{\beta_i\}_{i=1}^{L}$ を増加列とすると、前向き過程の周辺分布は

\[
p_L(x_L\mid x_0) \longrightarrow \mathcal{N}(0,I) \qquad (L\to\infty)
\]


へと収束し、これがデータ $x_0$ に依存しない事前分布

\[
p_{\mathrm{prior}} := \mathcal{N}(0,I)
\]


の選択を動機づける。


\subsection*{2.2.2 逆デノイジング過程（学習可能なデコーダ）}


その核心において、DDPM の本質は、前向き拡散過程が課す制御された劣化を逆転する能力にある。純で構造化されていないノイズ $x_L\sim p_{\mathrm{prior}}$ から始め、このランダム性をステップごとに徐々にデノイズし、一貫した意味のあるデータサンプルが現れるまで進めることが目標である。この逆生成は図2.5に示されるマルコフ連鎖を通じて進む。

\translationfigure{fig-2-5.jpg}{図2.5: DDPM 逆（デノイジング）過程の図示。ノイズ $x_L\sim p_{\mathrm{prior}}$ から始め、モデルは $i=L,\ldots,1$ について $x_{i-1}\sim p(x_{i-1}\mid x_i)$ を順にサンプリングし、新たに生成されたデータ $x_0$ を得る。オラクル遷移 $p(x_{i-1}\mid x_i)$ は未知であるため、それを近似することが目標である。}


根本的な課題、そして DDPM 開発を導く中心的な問いは次である：

\begin{propositionbox}[問い 2.2.1]
とくに $x_i\sim p_i(x_i)$ の複雑な分布を考えるとき、これらの逆遷移核 $p(x_{i-1}\mid x_i)$ を正確に計算できるか、あるいは少なくとも効果的に近似できるか？
\end{propositionbox}


元の DDPM 論文が行うように、直ちにエビデンス下界（Evidence Lower Bound; ELBO）の数学的に込み入った導出に飛び込むのではなく（その詳細な議論は第2.2.5節で待つ）、条件付き確率を活用して扱いやすい定式化を得る、より直感的な視点から学習目的にアプローチする。

\textbf{概観：モデリングと学習目的。} 生成過程を可能にするため、未知の真の逆遷移核 $p(x_{i-1}\mid x_i)$ を近似することが目標である。学習可能なパラメトリックモデル $p_\phi(x_{i-1}\mid x_i)$ を導入し、期待 KL ダイバージェンスを最小化するように学習することでこれを達成する：

\[
\mathbb{E}_{p_i(x_i)}\Bigl[D_{\mathrm{KL}}\bigl(p(x_{i-1}\mid x_i)\Vert p_\phi(x_{i-1}\mid x_i)\bigr)\Bigr]. \quad (2.2.2)
\]


しかし、目標分布 $p(x_{i-1}\mid x_i)$ の直接計算は困難である。ベイズの定理により、次を評価する必要がある：

\[
p(x_{i-1}\mid x_i) = p(x_i\mid x_{i-1})\,\frac{p_{i-1}(x_{i-1})}{p_i(x_i)}.
\]


ここで周辺 $p_i(x_i)$ と $p_{i-1}(x_{i-1})$ は、未知のデータ分布 $p_{\mathrm{data}}$ 上の期待値であり、次で与えられる：

\[
p_i(x_i) = \int p_i(x_i\mid x_0)\, p_{\mathrm{data}}(x_0)\, dx_0,
\]


および同様に $p_{i-1}(x_{i-1})$ についても。$p_{\mathrm{data}}$ が未知であるため、これらの積分は閉形式評価をもたず、せいぜいサンプルから近似できるにすぎず、厳密な密度は実際には利用できない。

\translationfigure{fig-2-6.jpg}{図2.6: DDPM における条件付けトリック。（左）固定されたノイズ付きサンプル $x_i$ に対し、逆分布 $p(x_{i-1}\mid x_i)$ は $x_i$ を生じうるすべての可能なクリーンサンプル $x_0$ にわたって平均する。したがって複雑な混合となりうる。（右）しかし学習中は、$x_i$ を生成するのに用いたクリーンサンプル $x_0$ が既知である。このサンプルで条件付けると $p(x_{i-1}\mid x_i,x_0)$ が得られ、これは閉形式の平均と分散をもつガウスである。したがって DDPM はこの扱いやすい条件付き分布を単純な回帰ターゲットとして用いることができる。各学習例は一つのクリーンサンプルのみで条件付けるが、条件付き学習目的をすべてのサンプルにわたって平均すると、周辺逆分布を直接学習するのと同じ最適モデルが得られる；定理 2.2.1 を参照。こうして DDPM は、周辺混合がもたらす複雑で一般に扱い不能なターゲットへ直接回帰することなく、同じ最適逆予測を学習できる。}


\textbf{条件付けによる扱い不能性の克服。} DDPM における中心的洞察がこの扱い不能性を解消する：逆遷移をクリーンデータサンプル $x_0$ で条件付ける。この繊細だが強力なステップが、扱い不能な核を数学的に扱いやすいものへと変換する：

\[
p(x_{i-1}\mid x_i,x_0) = p(x_i\mid x_{i-1})\,\frac{p(x_{i-1}\mid x_0)}{p_i(x_i\mid x_0)}.
\]


この扱いやすさは、前向き過程の二つの鍵性質から生じる：そのマルコフ性、すなわち $p(x_i\mid x_{i-1},x_0)=p(x_i\mid x_{i-1})$、および関与するすべての分布のガウス性である。その結果、$p(x_{i-1}\mid x_i,x_0)$ 自体がガウスであり、閉形式表現を認める（式 (2.2.4) で見る）。これを図2.6に可視化する。決定的に、この優雅な条件付け戦略により、式 (2.2.2) の一見扱い不能な周辺 KL ダイバージェンスと機能的に等価な扱いやすい目的を導出できる。

\begin{propositionbox}[定理 2.2.1: 周辺と条件付き KL 最小化の等価性]
次の等式が成り立つ：

\[
\begin{aligned}
&\mathbb{E}_{p_i(x_i)}\Bigl[D_{\mathrm{KL}}\bigl(p(x_{i-1}\mid x_i)\Vert p_\phi(x_{i-1}\mid x_i)\bigr)\Bigr] \\
&= \mathbb{E}_{p_{\mathrm{data}}(x_0)}\mathbb{E}_{p_i(x_i\mid x_0)}\Bigl[D_{\mathrm{KL}}\bigl(p(x_{i-1}\mid x_i,x_0)\Vert p_\phi(x_{i-1}\mid x_i)\bigr)\Bigr] + C. \quad (2.2.3)
\end{aligned}
\]

ここで $C$ は $\phi$ に依存しない定数である。さらに、式 (2.2.3) の最小化子は

\[
p^*(x_{i-1}\mid x_i) = \mathbb{E}_{p(x_0\mid x_i)}\bigl[p(x_{i-1}\mid x_i,x_0)\bigr] = p(x_{i-1}\mid x_i), \qquad x_i\sim p_i
\]

を満たす。
\end{propositionbox}


\textbf{定理の証明。}

証明は、定義を展開し、確率の連鎖律を適用し、対数恒等式を用いて、期待条件付き KL ダイバージェンスと周辺 KL ダイバージェンスの和へ分解することで、KL ダイバージェンス期待値を書き換える。完全な導出は第 D.1.1 節にある。■

この代替視点——扱いやすい目的を得るための条件付け——が DDPM の基礎をなし、他の影響力ある拡散モデルとの深い共通性を明らかにする。これを第3章および第5章で探求する。

それは強力な等価性を明らかにする：周辺分布間の KL ダイバージェンスの最小化は、特定の条件付き分布間の KL ダイバージェンスの最小化と数学的に同一である。後者の定式化は例外的に有用である。なぜなら、重要な条件付き分布 $p(x_{i-1}\mid x_i,x)$ が都合のよい閉形式表現をもつからである：

\begin{propositionbox}[補題 2.2.2: 逆条件付き遷移核]
$p(x_{i-1}\mid x_i,x)$ は次の閉形式表現をもつガウスである：

\[
p(x_{i-1}\mid x_i,x) = \mathcal{N}\bigl(x_{i-1};\, \mu(x_i,x,i),\, \sigma^2(i)I\bigr),
\]

ここで

\[
\mu(x_i,x,i) := \frac{\bar{\alpha}_{i-1}\beta_i^2}{1-\bar{\alpha}_i^2}\, x + \frac{(1-\bar{\alpha}_{i-1}^2)\alpha_i}{1-\bar{\alpha}_i^2}\, x_i, \qquad
\sigma^2(i) := \frac{1-\bar{\alpha}_{i-1}^2}{1-\bar{\alpha}_i^2}\,\beta_i^2. \quad (2.2.4)
\]
\end{propositionbox}


後に補題 4.5.2 で、式 (2.2.1) に記述された DDPM ノイズ過程を超えて拡張する、より一般的な公式を提示する。

\subsection*{2.2.3 逆遷移核 $p_\phi(x_{i-1}\mid x_i)$ のモデリング}


定理 2.2.1 における勾配レベルの等価性と、補題 2.2.2 における逆条件付き $p(x_{i-1}\mid x_i,x)$ のガウス形を活用し、DDPM は各逆遷移 $p_\phi(x_{i-1}\mid x_i)$ がガウスであると仮定し、次のようにパラメータ化する：

\[
p_\phi(x_{i-1}\mid x_i) := \mathcal{N}\bigl(x_{i-1};\, \mu_\phi(x_i,i),\, \sigma^2(i)I\bigr), \quad (2.2.5)
\]


ここで $\mu_\phi(\cdot,i):\mathbb{R}^D\to\mathbb{R}^D$ は学習可能な平均関数であり、$\sigma^2(i)>0$ は式 (2.2.4) で定義された通り固定である。

時間ステップ $i$ にわたって平均し、データ $x_0\sim p_{\mathrm{data}}$ で条件付けた KL ダイバージェンスを、すべての層の分布を整合させるものとして次のように記す：

\[
\mathcal{L}_{\mathrm{diffusion}}(x_0;\phi) := \sum_{i=2}^{L}\mathbb{E}_{p_i(x_i\mid x_0)}\Bigl[D_{\mathrm{KL}}\bigl(p(x_{i-1}\mid x_i,x_0)\Vert p_\phi(x_{i-1}\mid x_i)\bigr)\Bigr]. \quad (2.2.6)
\]


両分布のガウス形と式 (2.2.5) で定義されたパラメータ化のおかげで、目的は閉形式表現を認め、次のように単純化できる：

\[
\mathcal{L}_{\mathrm{diffusion}}(x_0;\phi) = \sum_{i=2}^{L}\frac{1}{2\sigma^2(i)}\,\mathbb{E}_{p_i(x_i\mid x_0)}\bigl\|\mu_\phi(x_i,i)-\mu(x_i,x_0,i)\bigr\|_2^2 + C, \quad (2.2.7)
\]


ここで $C$ は $\phi$ に依存しない定数である。データ分布にわたって平均し、定数 $C$（最適化に影響しない）を省略すると、最終的な DDPM 学習目的は

\[
\mathcal{L}_{\mathrm{DDPM}}(\phi) := \sum_{i=2}^{L}\frac{1}{2\sigma^2(i)}\,\mathbb{E}_{x_0}\mathbb{E}_{p_i(x_i\mid x_0)}\bigl\|\mu_\phi(x_i,i)-\mu(x_i,x_0,i)\bigr\|_2^2 \quad (2.2.8)
\]


であり、ここで $x_0\sim p_{\mathrm{data}}$ である。

前向き摂動核 $p_i(x_i\mid x_0)$ が学習目的において重要な計算的役割を果たすことに注意する。それは、クリーンサンプル $x_0$ から $i$ 回の前向きステップ後に得られるノイズ状態 $x_i$ の分布を与える。各前向き遷移が線形かつガウスであるため、その合成も閉形式のガウス分布をもつ；式 (2.2.1) 周辺の議論を参照。したがって、時間ステップ $i$ を選んだ後、$x_1,\ldots,x_{i-1}$ の中間状態をまず生成することなく、$x_0$ から直接 $x_i\sim p_i(x_i\mid x_0)$ をサンプリングできる。こうして各損失項は、クリーンサンプル、選ばれた時間ステップ、および直接劣化したサンプルのみを用いて推定でき、確率的学習を計算上効率的にする。

\subsection*{2.2.4 予測と損失の実践的選択}


**$\epsilon$-予測。** 典型的な DDPM 実装では、学習は式 (2.2.8) の平均予測パラメータ化に基づく元の損失を直接用いて行われない。代わりに、$\epsilon$-予測（ノイズ予測）定式化として知られる等価な再パラメータ化が一般に採用される。

DDPM 前向き過程では、ノイズレベル $i$ におけるノイズ付きサンプル $x_i\sim p(x_i\mid x_0)$ が次で生成されることを思い出そう：

\[
x_i = \bar{\alpha}_i x_0 + \sqrt{1-\bar{\alpha}_i^2}\,\epsilon, \qquad x_0\sim p_{\mathrm{data}},\quad \epsilon\sim\mathcal{N}(0,I). \quad (2.2.9)
\]


この表現を用いると、式 (2.2.4) の逆平均 $\mu(x_i,x_0,i)$ は次のように書き換えられる：

\[
\mu(x_i,x_0,i) = \frac{1}{\alpha_i}\Biggl(x_i - \frac{1-\alpha_i^2}{\sqrt{1-\bar{\alpha}_i^2}}\,\epsilon\Biggr).
\]


これは、ノイズを直接予測するニューラルネットワーク $\epsilon_\phi(x_i,i)$ を用いたモデル平均 $\mu_\phi$ のパラメータ化を動機づける：

\[
\mu_\phi(x_i,i) = \frac{1}{\alpha_i}\Biggl(x_i - \frac{1-\alpha_i^2}{\sqrt{1-\bar{\alpha}_i^2}}\,\underbrace{\epsilon_\phi(x_i,i)}_{\epsilon\text{-予測}}\Biggr).
\]


これを元の損失に代入すると、予測ノイズと真のノイズの間の二乗 $\ell_2$ 誤差が得られる：

\[
\bigl\|\mu_\phi(x_i,i)-\mu(x_i,x_0,i)\bigr\|_2^2 \propto \bigl\|\epsilon_\phi(x_i,i)-\epsilon\bigr\|_2^2,
\]


これは $i$ に依存する重み因子を除いてである。直感的に、モデルは「ノイズ探偵」として振る舞い、前向き過程の各ステップで加えられたランダムノイズを推定する。この推定を劣化したサンプルから引くとクリーンな元に近づき、このステップを繰り返すことで純ノイズからデータが再構成される。

**$\epsilon$-予測による単純化損失。** 実際には、この表現は重み項を省略することでさらに単純化され、広く用いられる DDPM 学習損失が得られる：

\[
\mathcal{L}_{\mathrm{simple}}(\phi) := \mathbb{E}_{i}\mathbb{E}_{x\sim p_{\mathrm{data}}(x)}\mathbb{E}_{\epsilon\sim\mathcal{N}(0,I)}\bigl\|\epsilon_\phi(x_i,i)-\epsilon\bigr\|_2^2, \quad (2.2.10)
\]


ここで $x_i = \bar{\alpha}_i x_0 + \sqrt{1-\bar{\alpha}_i^2}\,\epsilon$、$x_0\sim p_{\mathrm{data}}$ である。ターゲットノイズはすべての時間ステップ $t$ で単位分散をもつため、式 (2.2.10) の $\ell_2$ 損失はすべての $t$ にわたって一貫したスケールを保つ。これはターゲットの消失や爆発を防ぎ、明示的な損失重み付けの必要をなくす。

重要なことに、$\mathcal{L}_{\mathrm{DDPM}}$ と $\mathcal{L}_{\mathrm{simple}}$ は同じ最適解 $\epsilon^*$ を共有する。これは式 (2.2.10) が本質的に最小二乗問題に帰着するからである（命題 4.3.1 または命題 6.3.1 と同様に示される）\footnote{ELBO 由来の重みを $\mathcal{L}_{\mathrm{DDPM}}$ から除去すると学習目的は変わるが、母集団最適予測 $\epsilon^*(x_i,i)=\mathbb{E}[\epsilon\mid x_i]$ は変わらない。したがって重み付きと単純化された DDPM 損失は同じ理想予測器をもつが、学習中に異なるノイズレベルを強調する。実際には、この重み付けはなお最適化と性能に影響しうる。}：

\[
\epsilon^*(x_i,i) = \mathbb{E}[\epsilon\mid x_i], \qquad x_i\sim p_i.
\]


直感的に、$\epsilon$-予測ネットワーク $\epsilon_\phi(x_i,i)$ は、$x_i$ を生じるために前向き過程が加えたノイズを推定する。最適性において、この推定は真のノイズの条件付き期待値と一致する。$x_i$ が元のクリーンサンプルを一意に決定しなくてもである。

**もう一つの等価なパラメータ化：$x$-予測。** 式 (2.2.4) は、$x$-予測（クリーン予測）として知られる、別の等価なパラメータ化を動機づける。ここではニューラルネットワーク $x_\phi(x_i,i)$ が、ノイズレベル $i$ における与えられたノイズ付き入力 $x_i\sim p_i(x_i)$ からクリーン（デノイズされた）サンプルを予測するよう学習される。逆平均表現における真のクリーンサンプル $x_0$ を $x_\phi(x_i,i)$ で置き換えると、次のモデルが得られる：

\[
\mu_\phi(x_i,i) = \frac{\bar{\alpha}_{i-1}\beta_i^2}{1-\bar{\alpha}_i^2}\, x_\phi(x_i,i) + \frac{(1-\bar{\alpha}_{i-1}^2)\alpha_i}{1-\bar{\alpha}_i^2}\, x_i.
\]


$\epsilon$-予測定式化と同様、学習目的は次のように表せる：

\[
\bigl\|\mu_\phi(x_i,i)-\mu(x_i,x_0,i)\bigr\|_2^2 \propto \bigl\|x_\phi(x_i,i)-x_0\bigr\|_2^2, \qquad x_0\sim p_{\mathrm{data}},
\]


ここでモデルはノイズ付き版 $x_i$ から元データサンプル $x_0$ を予測するよう学習される。この等価性により、式 (2.2.8) の平均整合損失は、ある重み関数 $\omega_i$ に対して

\[
\mathbb{E}_{i}\mathbb{E}_{x_0\sim p_{\mathrm{data}}}\mathbb{E}_{\epsilon\sim\mathcal{N}(0,I)}\omega_i\bigl\|x_\phi(x_i,i)-x_0\bigr\|_2^2
\]


に帰着する。これは最小二乗問題であるため、最適解は（命題 4.3.1 または命題 6.3.1 を参照）

\[
x^*(x_i,i) = \mathbb{E}[x_0\mid x_i], \qquad x_i\sim p_i \quad (2.2.11)
\]


で与えられる。すなわち、モデルは時間ステップ $i$ におけるノイズ観測 $x_i$ が与えられたときの期待クリーンデータを予測すべきである。

$x$-予測と $\epsilon$-予測のパラメータ化は数学的に等価であり、前向き過程を通じて結ばれる：

\[
x_i = \bar{\alpha}_i x_\phi(x_i,i) + \sqrt{1-\bar{\alpha}_i^2}\,\epsilon_\phi(x_i,i). \quad (2.2.12)
\]


すなわち、クリーンサンプル $x_\phi(x_i,i)$ を予測しても、ノイズ $\epsilon_\phi(x_i,i)$ を予測してもよく、それらの組み合わせが前向きノイズ過程の下で $x_i$ を再現する。


\subsection*{2.2.5 DDPM の ELBO}


逆遷移を式 (2.2.5) のように定義すると、DDPM における同時生成分布の定義が次のように導かれる：

\[
p_\phi(x_0,x_{1:L}) := p_\phi(x_0\mid x_1)\, p_\phi(x_1\mid x_2)\cdots p_\phi(x_{L-1}\mid x_L)\, p_{\mathrm{prior}}(x_L),
\]


そしてデータの周辺生成モデルは次で与えられる：

\[
p_\phi(x_0) := \int p_\phi(x_0,x_{1:L})\, dx_{1:L}.
\]


実際、式 (2.2.6) による DDPM 学習は、最尤推定（式 (1.1.2)）に厳密に根拠づけられる。具体的には、その目的は第2.1節の VAE および HVAE におけるものと同様の ELBO をなし、対数密度の下界として機能する：

\begin{propositionbox}[定理 2.2.3: DDPM の ELBO]
\[
\begin{aligned}
-\log p_\phi(x_0) &\le -\mathcal{L}_{\mathrm{ELBO}}(x_0;\phi) \\
&:= \mathcal{L}_{\mathrm{prior}}(x_0) + \mathcal{L}_{\mathrm{recon.}}(x_0;\phi) + \mathcal{L}_{\mathrm{diffusion}}(x_0;\phi). \quad (2.2.13)
\end{aligned}
\]

ここで各損失成分は次で定義される：

\[
\begin{aligned}
\mathcal{L}_{\mathrm{prior}}(x_0) &:= D_{\mathrm{KL}}\bigl(p(x_L\mid x_0)\Vert p_{\mathrm{prior}}(x_L)\bigr), \\
\mathcal{L}_{\mathrm{recon.}}(x_0;\phi) &:= \mathbb{E}_{p(x_1\mid x_0)}\bigl[-\log p_\phi(x_0\mid x_1)\bigr], \\
\mathcal{L}_{\mathrm{diffusion}}(x_0;\phi) &= \sum_{i=2}^{L}\mathbb{E}_{p_i(x_i\mid x_0)}\Bigl[D_{\mathrm{KL}}\bigl(p(x_{i-1}\mid x_i,x_0)\Vert p_\phi(x_{i-1}\mid x_i)\bigr)\Bigr].
\end{aligned}
\]
\end{propositionbox}


\textbf{定理の証明。}
導出は、HVAE/VAE ELBO（式 (2.1.3)）と同様に Jensen の不等式を適用し、さらなる単純化を行う。詳細な証明は第 D.1.2 節に委ねる。■

ELBO $\mathcal{L}_{\mathrm{ELBO}}$ は三つの項からなる：

\begin{itemize}
\item $\mathcal{L}_{\mathrm{prior}}$ は、$p(\cdot\mid x_0)\approx p_{\mathrm{prior}}(\cdot)$ となるようノイズスケジュール $\{\beta_i\}$ を選ぶことで無視できる程度にできる。
\item $\mathcal{L}_{\mathrm{recon.}}$ については、モンテカルロ推定を用いて近似・最適化できる；実践的実装については (Ho et al., 2020; Kingma et al., 2021) を参照。
\item $\mathcal{L}_{\mathrm{diffusion}}$（式 (2.2.6) 参照）は、すべてのステップ $i$ で逆条件付き $p_\phi(x_{i-1}\mid x_i)$ を $p(x_{i-1}\mid x_i)$ に整合させる。
\end{itemize}


ELBO 目的 $\mathcal{L}_{\mathrm{ELBO}}$ は、潜在 $z=x_{1:L}$ をもつデータ処理不等式のレンズを通じても解釈できる。式 (2.1.2) に示されるように：

\[
D_{\mathrm{KL}}\bigl(p_{\mathrm{data}}(x_0)\Vert p_\phi(x_0)\bigr) \le D_{\mathrm{KL}}\bigl(p(x_0,x_{1:L})\Vert p_\phi(x_0,x_{1:L})\bigr),
\]


ここで $p(x_0,x_{1:L}):=p_{\mathrm{data}}(x_0)\,p(x_1\mid x_0)\,p(x_2\mid x_1)\cdots p(x_L\mid x_{L-1})$ は前向き過程に沿った同時分布を表す。

\textbf{注。}
拡散の変分的見方は、階層的潜在変数の描像に特に都合よく適合する：「エンコーダ」は固定前向きノイズ連鎖であり、学習された推論ネットワークはない。それでも ELBO は、式 (2.2.13) の終端事前項および再構成項とともに、ステップごとの KL 項へと分解する。DDPM 学習をとりわけ扱いやすくするのは、線形–ガウス前向き過程の下で、これらのステップごとの KL 項が、閉形式の条件付きターゲットをもつ単純なデノイジング回帰問題として書き換えられることである。したがって鍵となる単純化は、レベルごとの KL 項の不在ではなく、前向き過程が固定であり、それらを効率的にサンプリングできる教師ありデノイジング目的へと変えるという事実である。

\subsection*{2.2.6 サンプリング}


$\epsilon$-予測モデル $\epsilon_{\phi^\times}(x_i,i)$\footnote{記号「$\times$」は、モデルが学習済みで現在凍結されていることを示すために用いる。} を学習した後、図2.5に示されるように、パラメータ化された遷移 $p_{\phi^\times}(x_{i-1}\mid x_i)$ を用いて逐次的にサンプリングが行われる。

より具体的には、ランダムシード $x_L\sim p_{\mathrm{prior}}=\mathcal{N}(0,I)$ から始め、$i=L,L-1,\ldots,1$ について次の更新則に従って $p_{\phi^\times}(x_{i-1}\mid x_i)$ から再帰的にサンプリングする：

\[
x_{i-1} \leftarrow \underbrace{\frac{1}{\alpha_i}\Biggl(x_i - \frac{1-\alpha_i^2}{\sqrt{1-\bar{\alpha}_i^2}}\,\epsilon_{\phi^\times}(x_i,i)\Biggr)}_{\mu_{\phi^\times}(x_i,i)} + \sigma(i)\,\epsilon_i, \qquad \epsilon_i\sim\mathcal{N}(0,I). \quad (2.2.14)
\]


この「デノイジング」過程は、最終的なクリーン生成サンプルとして $x_0$ が得られるまで続く。


\textbf{DDPM サンプリングのもう一つの解釈。} DDPM サンプリングを解釈する有用な方法は、ノイズレベル構造が明示的になるよう更新則を書き換えることである。式 (2.2.12) から、予測ノイズは現在のサンプル $x_i$ と予測クリーンサンプル $x_{\phi^\times}(x_i,i)$ の観点で表せることを思い出そう：

\[
\epsilon_{\phi^\times}(x_i,i) = \frac{x_i - \bar{\alpha}_i x_{\phi^\times}(x_i,i)}{\sqrt{1-\bar{\alpha}_i^2}}.
\]


これを DDPM サンプリング則（式 (2.2.14)）に代入して整理すると、

\[
x_{i-1} \leftarrow \bar{\alpha}_{i-1}\, x_{\phi^\times}(x_i,i) + \underbrace{\sqrt{1-\bar{\alpha}_{i-1}^2-\tilde{\beta}_i^2}\,\epsilon_{\phi^\times}(x_i,i) + \tilde{\beta}_i\,\epsilon_i}_{\text{結合標準偏差}=\sqrt{1-\bar{\alpha}_{i-1}^2}}, \quad (2.2.15)
\]


が得られる。ここで $\tilde{\beta}_i^2 = \dfrac{(1-\bar{\alpha}_{i-1}^2)(1-\alpha_i^2)}{1-\bar{\alpha}_i^2}$ は DDPM 事後分散である\footnote{係数 $\sqrt{1-\bar{\alpha}_{i-1}^2-\tilde{\beta}_i^2}$ は $\dfrac{\alpha_i(1-\bar{\alpha}_{i-1}^2)}{\sqrt{1-\bar{\alpha}_i^2}}$ に単純化し、常に非負であるため、平方根は定義される。}。さらに、残差係数は

\[
\bigl(1-\bar{\alpha}_{i-1}^2-\tilde{\beta}_i^2\bigr) + \tilde{\beta}_i^2 = 1-\bar{\alpha}_{i-1}^2
\]


を満たすので、更新はレベル $i-1$ における前向き周辺と同じノイズレベル構造をもつ。これは式 (1.4.1) で導入された最初の操作的サンプリングテンプレートに戻る：クリーン部分を予測し、より低いノイズレベルでサンプルを書き換える。DDPM はこの書き換えを確率的に行い、$x_i$ から推定されたノイズ方向と新しいガウスノイズサンプルを組み合わせる。決定論的サンプラーは同じ基本的考えに従うが、新しいランダム性を除去し、残りのノイズ係数をそれに応じて調整する。拡散サンプラーをより系統的に研究する第9.2節で、この構成を再訪する。

式 (2.2.15) を前向き周辺 $x_{i-1}=\bar{\alpha}_{i-1}x_0+\sqrt{1-\bar{\alpha}_{i-1}^2}\,\bar{\epsilon}$ と比較せよ。構造は同一である：信号係数 $\bar{\alpha}_{i-1}$ とノイズ標準偏差 $\sqrt{1-\bar{\alpha}_{i-1}^2}$。ただし未知の $x_0$ がネットワーク予測 $x_{\phi^\times}(x_i,i)$ で置き換えられている。言い換えれば、各 DDPM ステップはクリーンデータを予測し、それを正確にノイズレベル $i-1$ へ再ノイズする。したがって DDPM サンプリングは二つの操作の反復として見ることができる：

\begin{enumerate}
\item \textbf{デノイズ。} 現在のノイズ付き入力 $x_i$ からクリーンデータ $x_{\phi^\times}(x_i,i)$ を推定する。
\item \textbf{再ノイズ。} 式 (2.2.15) を通じ、低減されたレベル $i-1$ でノイズを戻し加え、前向き過程のステップ $i-1$ におけるノイズレベルと一致するサンプル $x_{i-1}$ を生成する。
\end{enumerate}



しかし、$x_{\phi^\times}$ が最適デノイザ（すなわち条件付き期待値の最小化子；式 (2.2.11) を参照）として学習されていても、それは $x_i$ が与えられたときの平均クリーンサンプルしか予測できない。この限界は、とくに高ノイズレベルで、著しく劣化した入力から詳細構造を回復することが難しくなるぼやけた予測につながる。この視点から、拡散サンプリングは高ノイズから低ノイズへ移り、クリーン信号の推定を徐々に精錬する。初期ステップが大域構造を定め、後のステップが細かい詳細を加え、ノイズが除去されるにつれてサンプルはより現実的になる。

\translationfigure{fig-2-7.jpg}{図2.7: DDPM サンプリングのデノイズ後再ノイズの見方の図示。現在のノイズ付きサンプル $x_i$ から、モデルはまず背後のクリーン信号 $x_{\phi^\times}(x_i,i)$ を推定し、次いでこの推定を第 $(i-1)$ ノイズレベルへ再ノイズすることで、よりノイズの少ない点 $x_{i-1}$ をサンプリングする。}


\textbf{DDPM の遅いサンプリング速度。} DDPM（すなわち拡散モデル）におけるサンプリングは本質的に遅い。元の定式化と実装では、サンプル生成は典型的に約 1,000 のデノイジングステップを要する。これは、サンプリングが小さな精錬の長い系列を通じて進むからである：各ステップでモデルは現在のノイズ付きサンプルをわずかにクリーンなものへ更新し、次の更新は前のステップの結果から始めなければならない。その結果、標準 DDPM サンプリングは時間ステップ的で強く逐次的であり、良いサンプル品質は通常、逆拡散経路を密に追跡するために多くのそのようなステップを要する。

定理 2.2.1 は、表現力豊かな $p_\phi(x_{i-1}\mid x_i)$ が理論的に真の逆分布 $p(x_{i-1}\mid x_i)$ に整合しうることを示す。しかし実際には、$p_\phi(x_{i-1}\mid x_i)$ は典型的に $p(x_{i-1}\mid x_i)$ を近似するガウスとしてモデル化され、その表現力が制限される。

小さな前向きノイズスケール $\beta_i$ では、真の逆分布は近似的にガウスであり、正確な近似を可能にする。逆に、大きな $\beta_i$ は、単一ガウスが捉えられない多峰性または強い非ガウス性を誘発する。精度を保つため、DDPM は多くの小さな $\beta_i$ ステップを用い、各ステップが前のステップに依存しニューラルネットワーク評価 $\epsilon_{\phi^\times}(x_i,i)$ を要する逐次連鎖を形成する。これにより $O(L)$ の逐次パスが生じ、並列化を妨げ生成を遅くする。

後に第4章で、この固有のサンプリングボトルネックをより原理的に微分方程式問題として解釈し、生成を加速する連続時間数値戦略を動機づける様子を示す。
\end{document}
